\documentclass[11pt]{article}
\usepackage[a4paper,margin=27mm]{geometry}
\usepackage{amsmath,amssymb,mathtools}
\usepackage{fontspec,unicode-math}
\setmainfont{Latin Modern Roman}
\setmathfont{Latin Modern Math}
\usepackage{xurl,hyperref}
\hypersetup{colorlinks=true,urlcolor=blue,linkcolor=blue}
\newcommand{\sref}[2]{\hyperlink{source:#1:#2}{[S#2]}}
\newcommand{\cataloguescope}[1]{\paragraph{Mathematical question.} #1}
\begin{document}
% UnsolvedMath ID 8700011; source catalogue code AMR-086-0011
\setcounter{section}{8700010}
\section{Conjecture 2.3 — Erdős–Dressler}\label{8700011}
{\small\sffamily UnsolvedMath ID 8700011 · Graph Theory}
\subsection{Definitions and mathematical statement}

\paragraph{Shared notation.}$\mathbb N=\{1,2,\ldots\}$, and $\mathbb Z,\mathbb Q,\mathbb R,\mathbb C$ denote the integers, rationals, reals and complex numbers. Any other conventions are specified locally.


Use $R(a)=\prod_{p\mid a}p$ for the radical of a positive integer.
\paragraph{Statement recorded in the supplied source.}
For every $1\le a<b$ with $R(a)=R(b)$, must there be a prime $p$ satisfying $a<p<b$? \sref{8700011}{2}
\subsection{Short English statement}
Must two different positive integers with the same prime divisors have a prime strictly between them?
\subsection{Sources}
\begin{description}
\item[\hypertarget{source:8700011:1}{S1}] ULAM.AI and the UnsolvedMath Contributors, UnsolvedMath: A Curated Collection of Open Mathematics Problems, record 8700011 (AMR-086-0011). CC BY 4.0. \url{https://huggingface.co/datasets/ulamai/UnsolvedMath}.
\item[\hypertarget{source:8700011:2}{S2}] Michel Waldschmidt, Open Diophantine Problems, arXiv:math/0312440v2 (24 January 2004), Conjecture 2.3. \url{https://arxiv.org/abs/math/0312440}.
\end{description}
\label{end:8700011}
\end{document}
